\section{Déséquilibre de charge et splitting adaptatif}

\begin{figure}[H]
    \centering
    \begin{tikzpicture}
\begin{axis}[
    ybar stacked,
    width=\linewidth,
    height=7.0cm,
    ylabel={Temps (ms)},
    symbolic x coords={4w-1t,4w-2t,4w-4t,6w-2t,8w-1t},
    xtick=data,
    x tick label style={rotate=20, anchor=east},
    legend style={at={(0.5,-0.22)}, anchor=north, legend columns=2},
    enlarge x limits=0.15,
    grid=both,
    minor y tick num=1,
    axis line style={draw=black!60},
]
\addplot[fill=enstaBleuFonce] coordinates {
    (4w-1t, 74.338)
    (4w-2t, 45.457)
    (4w-4t, 34.728)
    (6w-2t, 35.606)
    (8w-1t, 46.872)
};
\addlegendentry{Autre calcul}

\addplot[fill=enstaBleuClair] coordinates {
    (4w-1t, 26.652)
    (4w-2t, 12.227)
    (4w-4t, 7.620)
    (6w-2t, 32.172)
    (8w-1t, 47.931)
};
\addlegendentry{Migration}

\addplot[fill=black!20] coordinates {
    (4w-1t, 0.493)
    (4w-2t, 0.488)
    (4w-4t, 0.537)
    (6w-2t, 0.627)
    (8w-1t, 0.658)
};
\addlegendentry{Halo}

\addplot[fill=black!40] coordinates {
    (4w-1t, 0.255)
    (4w-2t, 0.259)
    (4w-4t, 0.276)
    (6w-2t, 0.347)
    (8w-1t, 0.336)
};
\addlegendentry{Réduction}

\addplot[fill=black!65] coordinates {
    (4w-1t, 23.410)
    (4w-2t, 13.286)
    (4w-4t, 8.295)
    (6w-2t, 31.154)
    (8w-1t, 42.664)
};
\addlegendentry{Déséquilibre}
\end{axis}
\end{tikzpicture}

    \caption{Décomposition du temps critique sur plusieurs configurations stables de l'étape 3}
    \label{fig:stage3_breakdown}
\end{figure}

La figure précédente montre que le halo et la réduction globale restent très faibles. Le vrai problème n'est donc pas la communication de voisinage elle-même. Deux termes dominent beaucoup plus vite :
\begin{itemize}
    \item la migration des particules entre sous-domaines ;
    \item l'écart entre le worker moyen et le worker le plus lent.
\end{itemize}

Cette perte de balance vient de la densité non uniforme de la galaxie. Le centre contient beaucoup plus d'étoiles que la périphérie. Avec un découpage régulier en cellules, certains workers héritent d'une zone très dense et deviennent le chemin critique, tandis que d'autres finissent plus tôt et attendent.

Le cours suggère alors une idée naturelle : utiliser un \textit{splitting adaptatif}. Le principe serait de donner des sous-domaines plus petits aux zones denses et plus grands aux zones peu peuplées, afin de mieux répartir le nombre réel de particules et non seulement le nombre de cellules.

Cette solution améliorerait le balance de charge, mais elle introduirait un nouveau coût :
\begin{itemize}
    \item interfaces plus irrégulières ;
    \item gestion plus complexe des voisins ;
    \item migration plus difficile à maintenir ;
    \item halos moins réguliers et donc communication plus délicate.
\end{itemize}

Autrement dit, l'étape 3 valide bien la décomposition de domaine et les cellules fantômes, mais elle montre aussi pourquoi une grille régulière n'est pas suffisante pour obtenir une bonne scalabilité sur une galaxie très concentrée.
